\documentclass[12pt, border=8pt]{standalone}
\usepackage{tikz}
\usepackage{amsmath, amssymb, bm}
\usepackage{xcolor}
\usetikzlibrary{positioning, arrows.meta, calc, fit, backgrounds}

% ── Colour aliases ──────────────────────────────────────────────────────────
\definecolor{arrowblue}{RGB}{79,134,198}   % 4f86c6
\definecolor{cellgreen}{RGB}{92,184,92}    % 5cb85c
\definecolor{celltextgreen}{RGB}{45,122,45}% 2d7a2d
\definecolor{noteyellow}{RGB}{255,253,231} % fffde7
\definecolor{noteborder}{RGB}{240,173,78}  % f0ad4e
\definecolor{notetextbrown}{RGB}{125,92,0} % 7d5c00

\begin{document}
\pagestyle{empty}

  \textbf{RNN --- Unrolled Through Time}

\vspace{6pt}

\begin{tikzpicture}[
  >=Stealth,
  % RNN cell box
  rnnbox/.style={
    draw=cellgreen, fill=cellgreen!12, rounded corners=5pt,
    minimum width=3.4cm, minimum height=1.8cm,
    align=center, font=\small, line width=1.5pt
  },
  % gray init box
  initbox/.style={
    draw=gray!50, fill=gray!8, rounded corners=5pt,
    minimum width=1.6cm, minimum height=1.2cm,
    align=center, font=\small
  },
  % horizontal hidden-state arrow
  harrow/.style={->, arrowblue, line width=1.5pt},
  % vertical input arrow
  xarrow/.style={->, gray!70, thin}
]

% ─── h0 init box ─────────────────────────────────────────────────────────────
\node[initbox] (h0box) at (0, 0) {$h_0$\\\textcolor{gray!60}{\footnotesize(zeros)}};

% ─── RNN Cell t=1 ────────────────────────────────────────────────────────────
\node[rnnbox, right=1.2cm of h0box] (cell1) {
  \textcolor{celltextgreen}{\textbf{RNN Cell \; $t=1$}}\\[4pt]
  $h_1 = \tanh(W_h h_0 + W_x x_1)$
};

% ─── RNN Cell t=2 ────────────────────────────────────────────────────────────
\node[rnnbox, right=1.6cm of cell1] (cell2) {
  \textcolor{celltextgreen}{\textbf{RNN Cell \; $t=2$}}\\[4pt]
  $h_2 = \tanh(W_h h_1 + W_x x_2)$
};

% ─── RNN Cell t=3 ────────────────────────────────────────────────────────────
\node[rnnbox, right=1.6cm of cell2] (cell3) {
  \textcolor{celltextgreen}{\textbf{RNN Cell \; $t=3$}}\\[4pt]
  $h_3 = \tanh(\ldots)$
};

% ─── Horizontal h-state arrows ───────────────────────────────────────────────
% h0 box → cell1
\draw[harrow] (h0box.east) -- (cell1.west)
  node[midway, above, font=\itshape\small, text=arrowblue] {$h_0$};

% cell1 → cell2
\draw[harrow] (cell1.east) -- (cell2.west)
  node[midway, above, font=\itshape\small, text=arrowblue] {$h_1$};

% cell2 → cell3
\draw[harrow] (cell2.east) -- (cell3.west)
  node[midway, above, font=\itshape\small, text=arrowblue] {$h_2$};

% exit arrow after cell3 (red, continuation)
\draw[->, red!70!black, line width=1.5pt] (cell3.east) -- ++(0.6, 0);

% ─── x inputs from below ─────────────────────────────────────────────────────
\draw[xarrow] (cell1.south) ++(0,-0.7) -- (cell1.south)
  node[below=0.7cm, font=\itshape\small, text=gray!70] {$x_1$};

\draw[xarrow] (cell2.south) ++(0,-0.7) -- (cell2.south)
  node[below=0.7cm, font=\itshape\small, text=gray!70] {$x_2$};

\draw[xarrow] (cell3.south) ++(0,-0.7) -- (cell3.south)
  node[below=0.7cm, font=\itshape\small, text=gray!70] {$x_3$};

% ─── Shared-weights note box ─────────────────────────────────────────────────
\node[draw=noteborder, fill=noteyellow, rounded corners=3pt,
      minimum width=12.5cm, minimum height=0.7cm,
      align=center, font=\small, text=notetextbrown,
      below=2.2cm of cell2] (note) {
  $\mathbf{\text{⚠}}$\quad
  $W_h$,\; $W_x$,\; $b$ are \textbf{shared} across all time steps --- not separate per step
};

\end{tikzpicture}

\end{document}
